\minitopic{运气不是单一事物}“幸运地知道”并不矛盾：研究者偶然翻到关键档案，随后认真核验并获得知识。构成威胁的是\term{认识运气}{epistemic luck}：在与实际情形非常接近的可能情况下，主体以同样方式相信 $p$，但 $p$ 为假。发现证据的运气与信念为真的运气应当区分\parencite{pritchard2005}。 反事实理论借助\term{可能世界}{possible worlds}表达这种邻近性。可能世界不是必须相信的平行宇宙，而是谈论“事情本可怎样”的模型。哪些世界算近，取决于保持哪些实际条件不变；这正是理论争议所在。

\minitopic{敏感性：信念是否追踪真值}诺齐克提出四项条件：$p$为真；$S$相信$p$；若$p$不为真，$S$不会相信$p$；若$p$仍为真但情形略有变化，$S$仍会相信$p$\parencite{nozick1981}。第三项通常称为\term{敏感性}{sensitivity}：
\[
\neg p \Box\!\!\to \neg B_Sp .
\]
它要求信念随真值变化。温度计若在温度不同的近邻情形仍显示相同读数，就不追踪温度。 敏感性自然处理许多Gettier案例，却与闭合原则冲突。我可能通过正常视觉知道自己有手，但在最近的“我没有手”世界——缸中脑世界——仍会相信自己有手，因此“我不是缸中脑”不敏感。追踪理论可能允许我知道有手，却不知道由它演绎出的反怀疑命题。

\minitopic{安全性：是否容易出错}\term{安全性}{safety}把方向倒过来：在主体以相同基础相信$p$的近邻世界中，$p$不会为假。粗略地说，安全信念不容易错。它比敏感性更容易保留闭合，也能解释谷仓立面：亨利若稍看旁边就会以同样方式形成假信念，所以实际真信念不安全。 安全性仍面对三项选择。第一，邻近程度如何确定？第二，“相同基础”指同一心理状态、方法类型还是具体能力？第三，知识是否允许极小错误风险？把安全要求设得过强会复活怀疑论，过弱则放过运气。

\minitopic{德性：成功是否归功于能力}德性认识论把知识比作技艺表现。一个射击行动若\textbf{准确}、\textbf{熟练}，并且因为熟练而准确，就具有“恰当性”。相应地，信念需要为真，出自认知能力，并且其为真归功于能力\parencite{sosa2007}。这里的\term{认知德性}{epistemic virtue}既可指可靠的视知觉、记忆等能力，也可指开放、认真、谦逊等人格品质\parencite{zagzebski1996}。 能力解释比单纯因果关系更细：开篇射手的命中经过干预抵消，不适当地归功于能力。然而“归功”存在程度。友善的环境、仪器和他人帮助往往也是成功条件。如果要求主体单独拥有全部功劳，团队科学和证言知识就会被不合理排除。

\minitopic{知识优先的转向}也许不断补充条件的项目方向错了。知识优先纲领把知识视为不可由信念、真、确证等更基础概念非循环分析的状态，并用知识解释证据、断言和行动\parencite{williamson2000}。它不是说知识神秘不可研究，而是改变解释方向。例如“主体的证据就是其知识”把知识放在证据理论起点。 知识优先避免了定义竞赛，却仍须回答哪些案例是知识以及为什么。如果判断仍暗中依赖安全性或能力，理论分歧可能只是重新安排了概念顺序。

\begin{comparison}
儒家工夫论常把“知”置于长期形成的判断和行动能力中。它与德性认识论都反对把认识主体当作被动容器。但“德性”在儒家文本中具有伦理、关系和礼仪维度，不能直接等同于Sosa意义上的可靠认知能力。可比之处是对稳定能力的重视，差异在于能力服务的生活目标。
\end{comparison}

\minitopic{怎样选择反事实世界}反事实判断不是把一切可能情形平均计算。评估“如果杯子掉下去就会碎”，我们通常保持杯子材质、重力和地面不变，只改变它是否掉落；不会优先考虑魔法介入的遥远世界。同样，评估信念安全性时要保持形成方法和背景能力相对稳定。不过何者属于背景，正可能受争议：谷仓地区的假立面是方法外部的偶然环境，还是视觉辨认任务的一部分？ 世界接近性不能完全由物理相似决定。一个彩票号码只改一位在物理上很小，却改变中奖真值；一个遥远星系变化很大，却与当前信念无关。认识论接近性围绕信念形成条件组织，包括主体会注意什么、信息渠道如何响应事实、错误是否由同一缺陷产生。

\minitopic{能力的层级}说主体“有视力”太粗。能力总是在领域和条件下表现：白天辨认普通物体、昏暗中辨认颜色、在假立面环境辨认谷仓是不同能力—环境组合。可靠主义的过程类型问题在德性理论中以能力划分重新出现。划得太宽会把不可靠情形平均进去，划得太窄则可把每次成功描述成“在恰好这次条件下辨真”的完美能力。 一种约束是使用解释上自然、可跨案例重复并在主体行为中稳定体现的能力类型。训练、校准和错误模式可提供证据。能力也可分为一阶与二阶：一阶视觉产生对象信念，二阶监控判断光线是否适合继续依赖视觉。反思知识不必要求对全部机制有理论知识，却可能要求适当的二阶视角。

\minitopic{案例工作坊：导航与环境变化}司机依赖过去一直可靠的导航。道路昨夜封闭，但系统地图未更新；临时工作人员错误移走封路牌，使旧路线今天恰好仍可通行。司机相信路线畅通并成功抵达。该信念敏感吗？在最近的道路真正封闭情形，系统仍会说畅通，所以不敏感。安全吗？在许多近邻情形中工作人员没有移牌，信念为假，所以不安全。成功是否归功于导航能力？至多部分归功，关键真值来自偶然干预。 若系统实时接收道路传感器并在封闭时改道，信念开始追踪事实；若司机同时看见道路开放，形成基础也改变。这个案例提醒我们，认识能力经常分布在人、设备和维护制度之间。评价不能只问“司机可靠吗”或“软件准确率多高”，还要问整条信息链能否随目标事实变化。
